\clearpage
\thispagestyle{empty}
\pdfbookmark[0]{Allgemeine Lehrsätze. Ausgabe und Quellen}{lehrsaetze-source-notice}
\begin{center}\Large\textbf{Allgemeine Lehrsätze über Anziehungs- und Abstoßungskräfte}\end{center}

KI-Übertragung des deutschen Textes aus Carl Friedrich Gauss, \textit{Werke}, Band V (Göttingen, 1867), S.~195--242, in modernes Deutsch mit KI-erstelltem Satz. Die Ausgabe enthält das vollständige Titelblatt mit der historischen Verlagsangabe, die wiederholte Überschrift, alle 37 Artikel und die Anmerkung von Gauß zu Artikel 12. Die Abhandlung erschien ursprünglich in den \textit{Resultaten aus den Beobachtungen des magnetischen Vereins im Jahre 1839}, herausgegeben von Gauß und Weber (Leipzig, 1840).

Zum Vergleich dienten die englische Fassung in Richard Taylor (Hrsg.), \textit{Scientific Memoirs}, Band III, Teil X, Beitrag VII (1843), S.~153--196, sowie der von A.~K.~T.~Assis herausgegebene und kommentierte Anfangsauszug vom Oktober 2024. Für die heutige deutsche Terminologie wurde außerdem Alexander Mirlins \textit{Skript Elektrodynamik}, Kapitel I, \textit{Elektrostatik}, am KIT herangezogen.

\noindent\href{https://gdz.sub.uni-goettingen.de/id/PPN236006339}{Deutsche Vorlage (1867)}\quad
\href{https://archive.org/details/scientificmemoir03memo/page/153/mode/1up}{Englische Fassung (1843)}\quad
\href{https://www.ifi.unicamp.br/~assis/Gauss-1840.pdf}{Anfangsauszug von Assis (2024)}\quad
\href{https://www.tkm.kit.edu/downloads/ws2012_theoc/Kapitel-I-Elektrostatik.pdf}{Mirlin, Elektrostatik}
\clearpage
